% Folder 144, batch 07 — archivists' pages 121–140 (PDF pages 1–20 of the batch file).
% Pass: Opus 5.5 (claude-opus-5-5), 2026-10-03 — first pass, unchecked against the pages by a human.
%
% Pages 121-123 are the end of a letter of P. Deligne (begun in batch 6):
% third-party correspondence, summarised in a note and not transcribed
% (RIGHTS.md).
% Skipped (no mathematics): 124, 127, 132 (blank sheets), 126 (an envelope),
% 131, 134, 136, 138 (typed administrative notices, nothing in his hand).
% The pencilled numbers on the leaves read one higher than the \page numbers
% used here (the leaf at PDF page 1 carries "122"); \page follows the
% batch layout given by the parent run (PDF page 1 = 121).
\documentclass[11pt,a4paper]{article}
\input{../preamble/grothendieck}

\folder{144}
\batch{7}
\pages{121}{140}
\foldertitle{[Suite autour de Teichmüller dont] Teichmülleries, 1981-1982 : notes manuscrites (1981-1983, s.d.), lettres (1981, s.d.).}
\dating{1981-1983}
\watermark{Édition de démonstration}

\begin{document}

\section{Lettre : multiplication complexe et action de Galois}

\page{121}
\note{Les pp.~121 à 123 sont la fin de la seconde lettre de Pierre Deligne à Grothendieck, commencée au lot~6 (pp.~119--120), signée « Pierre » p.~123 : énoncé du théorème de Shimura--Taniyama et action de $\mathrm{Gal}(\overline{\mathbb{Q}}/K)$ sur les classes de courbes elliptiques à multiplication complexe par le groupe des classes d'idéaux. Correspondance d'un tiers, non transcrite (RIGHTS.md).}

\page{122}

\page{123}

\section{Groupes $\mathfrak{S}$, revêtements $\mathcal{M}$ et indices}

\page{125}
\note{Feuille de papier kraft ; encre et crayon, sans phrases. On note $\mathfrak{S}$ la lettre bouclée qu'il emploie aussi bien pour le groupe symétrique $\mathfrak{S}_3$ que pour les groupes $\mathfrak{S}_{0,3}$, $\mathfrak{S}_{1,1}$, $\mathfrak{S}^+$ des pages suivantes ; $\mathcal{M}$ (souvent souligné) pour ses multiplicités modulaires.}

Diagramme d'indices (en haut à gauche) :

\begin{tikzcd}[column sep=small, row sep=small]
\mathcal{M}^!_0 \arrow[r, no head, "2"] \arrow[d, no head, "4"'] & \underline{\mathcal{M}}^! \arrow[r, no head, "6"] \arrow[d, no head, "4"'] & \mathcal{M} \arrow[d, no head, "4"] \\
\underline{\mathcal{M}}^{!\circ}_0 \arrow[r, no head, "2"] \arrow[d, no head] & \underline{\mathcal{M}}^{!\circ} \arrow[r, no head, "6"] \arrow[d, no head] & \underline{\mathcal{M}}^{\circ} \\
1 \arrow[r, no head, "2"] & \mu &
\end{tikzcd}

\note{Le « 6 » de la deuxième ligne est surchargé (peut-être biffé) ; la lecture des exposants $^{!\circ}$ est incertaine.}

\[
  \mathcal{M}^{!\circ}_0 \simeq \mathcal{M}(0)^{\circ}\cdot\Pi_0
\]
avec, sous les deux facteurs, $\cap\,4$ et $\cap\,12$, puis $\mathcal{M}(0)\cdot\mathfrak{S}^+$.

\drawing{à droite, sous l'inscription $\mathcal{M}''^!_0$, un réseau en perspective (double pyramide) dont les sommets sont $\mathcal{M}^!_0$, $\underline{\mathcal{M}}^!$, $\underline{\mathcal{M}}$, $\underline{\mathcal{M}}'$, $\mathcal{M}'''$, $\underline{\mathcal{M}}''$, $\mathcal{M}'^!_0$, $\mathcal{M}'''^!_0$, $\mathcal{M}''^!_0$, $\mathcal{M}'^!$, $\mathcal{M}'''^!$, $\mathcal{M}''^!$, $\mathcal{M}^{!\circ}$, $\mathcal{M}^{\circ}$, $\underline{\mathcal{M}}^{!\circ}_0$ ; les arêtes portent les indices 2, 6, 12, et des arêtes pointillées verticales relient les sommets $'''$ aux sommets voisins. Dessous, au crayon, l'esquisse d'un cube $\mathcal{M}^!_0$, $\mathcal{M}^!$, $\mathcal{M}$, $\mathcal{M}^{!\circ}_0$, $\mathcal{M}^{!\circ}$ et d'autres polyèdres non étiquetés.}

\[
  D_\rho,\ D_\sigma,
\]
\[
  \mathcal{M} \longrightarrow \mathfrak{S}/\Pi_0 \times \mu \qquad (\varphi,\varepsilon_3)
\]
avec, au-dessus, une flèche verticale $\mathcal{M}^{!\circ}_0 \to$ marquée $\mu\times 1$, et
\[
  \mathcal{M} \longrightarrow \mathfrak{S}_3 \times\mu\times\mu \longrightarrow \mathfrak{S}_3\times\mu\times\mu,
  \qquad \mathfrak{S}_3\times\mu .
\]
\[
  \underline{\mathcal{M}}/\underline{\mathcal{M}}^{!\circ} \simeq \mathfrak{S}_3\times\mu\times\mu,
  \qquad
  \underline{\mathcal{M}}/\underline{\mathcal{M}}^{!\circ}_0 \longrightarrow \cdots
\]
\[
  1\to\mu\to\cdots\qquad
  1\to \mathrm{Sl}(2,\mathbb{Z}/4\mathbb{Z}) \to \mathrm{Gl}(2,\mathbb{Z}/4\mathbb{Z}) \to (\mathbb{Z}/4\mathbb{Z})^* \to 1
\]
sous $\mathrm{Sl}(2,\mathbb{Z}/4\mathbb{Z})$ : $48$ (\struck{48}), une flèche marquée $8$ vers
$6$ $\mathrm{Sl}(2,\mathbb{Z}/2\mathbb{Z})$, et \uncertain{$3\cdot 16$}.
\note{Le premier chiffre de « $3\cdot16$ » est écrit sur un 6.}

$\mathbb{F}_2$ --- \uncertain{$\mathbb{P}^1(\mathbb{F}_2)$}
\[
  \mathcal{M}''^!_\bullet \cap N^*_\rho = Z^{''!}_\rho
\]
\note{Dans le membre de droite, une astérisque surchargée sous les accents ; l'indice $\bullet$ de $\mathcal{M}''^!_\bullet$ peut être un $0$.}

\section{Le groupe $\mathfrak{S}_{1,1}$ et ses générateurs}

\page{128}
\[
  \mathfrak{S}_{1,1} \simeq \mathcal{T}_{1,2}^{\,!}
\]
\drawing{au crayon, un réseau carré de points (une douzaine de points pleins) sur lequel sont tracées des flèches horizontales et verticales marquées $\tilde u$, $\tilde v$, et des diagonales marquées $\tilde\rho$, $\tilde\sigma$ ; en haut à droite un cadre fléché.}

Au crayon :
\[
  \begin{pmatrix}-1 & 1\\ 0 & 1\end{pmatrix}
\]
$\rho($\ill{}$,\tfrac{1}{2}) = (\tfrac12, -\tfrac12)$
\note{Le premier argument de $\rho$ est biffé ; le dénominateur du second se lit mal.}
\[
  r_0\,u\,r_0 \qquad (u r_1)^{-1} = r_1 u^{-1} \qquad (u r_\infty)^{-1} = r_\infty u^{-1}
\]
\[
  (u r_\infty)^{-1} \qquad (v r_\infty)^{-1} = r_\infty v^{-1} \qquad r_1 v^{-1}
\]
À l'encre :
\[
  \tilde v^{-1}\,\tilde u\,\tilde v\,\tilde u^{-1} \qquad \tilde u\,\tilde v \qquad \tilde\rho
\]
\[
  (u,l)(u',l') = \bigl(uu',\, u(l')\,l\bigr)
\]
\note{Le dernier facteur est lu « $l$ » ; il pourrait s'agir d'un « $e$ ».}
\drawing{un petit triangle : $a \xrightarrow{\,l\,} u(a)$, $a\xrightarrow{\,l'\,} u'(a)$, puis $u(a)\to uu'(a)$ et $u'(a)\to uu'(a)$ ; une étiquette biffée.}
\[
  l : a\to ua,\qquad l' : a\to u'a, \qquad
  u'(l) : u'a \to u'ua,\qquad u(l') : ua \to uu'a
\]
\[
  a \xrightarrow{\ l\ } ua \longrightarrow uu'a
\]

\page{129}
\note{Notes au crayon et à l'encre en marge d'une affiche dactylographiée du séminaire de géométrie algébrique (USTL, novembre--décembre, avec un tampon « 30.10.80 ») ; seule l'écriture manuscrite est transcrite.}
\[
  (\gamma,l)(\gamma',l')(\gamma'',l'')
\]

\begin{tikzcd}[column sep=small]
a \arrow[r, "{l''}"] & \gamma''^{-1}(a) \arrow[r, "{\gamma''^{-1}(l')}"] & (\gamma'\gamma'')^{-1}(a) \arrow[d, "{(\gamma'\gamma'')^{-1}(l)}"] \\
 & & (\gamma\gamma'\gamma'')^{-1}(a)
\end{tikzcd}

\note{À l'encre, sur la même feuille, des mots en $\tilde u,\tilde v$ (avec exposants $\pm1$), en partie biffés de deux traits obliques, et deux lignes de calcul qu'on lit avec peine :}
\[
  \tilde v = [\tilde v,\tilde u^{-1}] \ldots
\]
\uncertain{$\tilde v = [\tilde v\,\tilde u]\,\tilde u^{-1} = (\tilde u\,\tilde v$} \ill{}

\page{130}
\[
  \tilde u,\ \tilde v \qquad [\tilde u,\tilde v] = l^{!\,-1}_0
\]
\[
  \tilde\rho,\ \tilde\sigma,\ \tilde\omega_0
\]
\[
  \tilde\rho^{\,3} = \tilde\omega_0,\qquad \tilde\sigma^{2} = \tilde\omega_0^{-1}
\]
\[
  \tilde\omega_0^{\,2} = l^!_0
\]
\note{Devant $l^{!\,-1}_0$ dans la première ligne, un signe biffé ; après $\tilde\omega_0^{\,2} = l^!_0$, une suite biffée terminée par « $=1$ ».}
\[
  \left.
  \begin{aligned}
    \tilde r_0(\tilde\rho) &= \tilde\rho^{-1}\\
    \tilde r_0(\tilde\sigma) &= \tilde\sigma^{-1}
  \end{aligned}
  \right\}
  \Longrightarrow\ r_0(\tilde\omega_0) = \tilde\omega_0^{-1}
\]
\[
  \text{déf}
  \begin{cases}
    \tilde r_1 = \tilde r_0\,\tilde\rho & \text{i.e. } \tilde\rho = \tilde r_0\,\tilde r_1\\
    \tilde r_\infty = \tilde\sigma\,\tilde r_0 & \text{i.e. } \tilde\sigma = \tilde r_\infty\,\tilde r_0
  \end{cases}
\]
\[
  \boxed{\ \tilde r_0^{\,2} = \tilde r_1^{\,2} = \tilde r_\infty^{\,2} = 1,\quad
  (\tilde r_0\tilde r_1)^3(\tilde r_\infty\tilde r_0)^2 = 1\ }
\]
En regard :
\[
  \mathfrak{S}_{1,1} \simeq \mathcal{T}_{1,2}^{\,!}
  \ \supset\ S\mathcal{T}_{1,1} \simeq \mathrm{Gl}(2,\mathbb{Z})^{\sim}
  = \mathrm{Norm}_{\mathfrak{S}_{1,1}}(L^!_0)
\]
($L^!_0$ engendré par $l^!_0$)

NB \uncertain{$\tilde r_1$}$\,\tilde r_\infty = (\tilde\sigma\tilde\rho)^{-1} = \tilde\varepsilon_0^{-1}$
où $\tilde\varepsilon_0 \overset{\text{déf}}{=} \tilde\sigma\tilde\rho$
\[
  \begin{aligned}
    \tilde\rho(\tilde u) &= (\tilde r_0\tilde r_1)(\tilde u) = \tilde v^{-1}\\
    \tilde\rho(\tilde v) &= (\tilde r_0\tilde r_1)(\tilde v) = \tilde v\,\tilde u\\
    \tilde\sigma(\tilde u) &= (\tilde r_\infty\tilde r_0)(\tilde u) = \tilde u(\tilde v) = \tilde u\,\tilde v\,\tilde u^{-1}\\
    \tilde\sigma(\tilde v) &= (\tilde r_\infty\tilde r_0)(\tilde v) = \tilde u^{-1}
  \end{aligned}
\]
\note{Dans chacune de ces quatre lignes, une expression biffée et illisible précède le second membre.}
\[
  \left.
  \begin{aligned}
    \tilde r_0(\tilde u) &= \tilde r_0\,\tilde u\,\tilde r_0 = \tilde v\\
    \tilde r_0(\tilde v) &= \tilde r_0\,\tilde v\,\tilde r_0 = \tilde u
  \end{aligned}
  \right\}
  \Longrightarrow\ \tilde r_0(l^!_0) = l^{!\,-1}_0
\]
\[
  \begin{aligned}
    \tilde r_\infty(\tilde u) &= \tilde r_\infty\,\tilde u\,\tilde r_\infty = \tilde u^{-1}\\
    \tilde r_\infty(\tilde v) &= \tilde r_\infty\,\tilde v\,\tilde r_\infty
    = \tilde v^{\,?}\,[\tilde v^{-1},\tilde u]\\
    &= \tilde u\,\tilde v\,\tilde u^{-1} = \tilde u(\tilde v)
  \end{aligned}
\]
\note{Dans le calcul de $\tilde r_\infty(\tilde v)$, deux expressions biffées ; l'exposant noté $?$ est illisible.}
$\Longrightarrow$ \uncertain{$\tilde r_\infty(l^!_0) = l^!_0$}
\note{Cette conclusion est ajoutée au bord droit, serrée ; l'indice de $\tilde r$ s'y lit plutôt « 1 ».}
\[
  \left.
  \begin{aligned}
    \tilde r_1(\tilde u) &= \tilde r_1\,\tilde u\,\tilde r_1 = \tilde u^{-1}\\
    \tilde r_1(\tilde v) &= \tilde r_1\,\tilde v\,\tilde r_1 = \tilde u\,\tilde v
  \end{aligned}
  \right\}
  \Longrightarrow\ \tilde r_1(l^!_0) = l^{!\,-1}_0
\]
D'où la présentation :
\[
  \Bigl\{\tilde r_0,\tilde r_1,\tilde r_\infty,\tilde u,\tilde v \ \Bigm|\
  \begin{array}{l}
    \tilde r_0^{\,2} = \tilde r_1^{\,2} = \tilde r_\infty^{\,2}
      = 1\\
    (\tilde r_0\tilde r_1)^{6}\ \bigl(= (\tilde r_\infty\tilde r_0)^4\bigr) = [\tilde u,\tilde v]\\
    \tilde r_0(\tilde u) = \tilde v,\ \tilde r_0(\tilde v) = \tilde u\\
    \tilde r_\infty(\tilde u) = \tilde u^{-1},\ \tilde r_\infty(\tilde v) = \tilde u(\tilde v)\ \bigl(= \tilde u\tilde v\tilde u^{-1}\bigr)\\
    r_1(\tilde u) = \tilde u^{-1};\ r_1(\tilde v) = \tilde u\,\tilde v
  \end{array}
  \Bigr\}
\]
\[
  = \Bigl\{\tilde r_0,\tilde r_1,\tilde r_\infty,\tilde\omega_0,\tilde u,\tilde v \ \Bigm|\
  \begin{array}{l}
    r_0^2 = r_1^2 = r_\infty^2 = 1,\ (\tilde r_0\tilde r_1)^3 = (\tilde r_\infty\tilde r_0)^2 = \tilde\omega_0\\
    \tilde\omega_0^{\,2} = [\tilde u,\tilde v]\\
    \text{relations de conjugaison (C)}
  \end{array}
  \Bigr\}
\]
\note{Dans la première relation, après $\tilde r_\infty^{\,2} =$, une ligne biffée (une première écriture de $(\tilde r_0\tilde r_1)^3(\tilde r_\infty\tilde r_0)^2 = 1$, recopiée au-dessus). L'exposant 6 de $(\tilde r_0\tilde r_1)$ est surchargé et douteux.}

\section{Le cas $(0,3)$ : positions relatives et diagrammes d'indices}

\page{133}
\[
  \begin{array}{ccccc}
    1 & & & & 1\\
    \uparrow & & & & \uparrow\\
    \mathfrak{S}_{0,3} & \longrightarrow & \mathcal{T}_{0,3} & = & \mathfrak{S}_{0,3}/\mathfrak{S}^{+!}_{0,3} \simeq \mathfrak{S}_3\times\mu_\tau\\
    \uparrow & & \uparrow & & \\
    \widetilde{\mathfrak{S}}^{+} & \longrightarrow & \widetilde{\mathcal{T}}_{0,3} & = & \cdots/(\mathfrak{S}^{!+}_{0,3})_0\\
    \uparrow & & \uparrow & & \\
    \mu & = & \mu & & \\
    \uparrow & & \uparrow & & \\
    1 & & 1 & &
  \end{array}
\]
\note{Le groupe de la colonne de gauche, lu $\widetilde{\mathfrak{S}}^{+}$, est surchargé et douteux ; le numérateur de $\widetilde{\mathcal{T}}_{0,3}$ est un symbole biffé illisible. À côté du $1$ du haut de la colonne de droite, un mot biffé. $(\mathfrak{S}^{!+}_{0,3})_0$ est souligné d'une accolade : « engendré par $\lambda_0\lambda_1,\lambda_\infty$ ($\lambda_i = -l_i$) ».}
\[
  \mathfrak{S} = \Bigl\{\rho,\sigma,\tau' \ \Bigm|\
  \begin{array}{l}
    \rho^6 = \sigma^4 = \tau'^2 = 1,\ \rho^3 = \sigma^2;\\
    \tau'(\rho) = \rho^{-1},\ \tau'(\sigma) = \sigma^{-1}
  \end{array}\Bigr\}
\]
\[
  = \Bigl\{\rho,\sigma,\tau \ \Bigm|\
  \begin{array}{l}
    \rho^6 = \sigma^4 = \tau^2 = 1,\ \rho^3 = \sigma^2\\
    \tau(\rho) = \cdots,\ \tau(\sigma) = \sigma^{-1}
  \end{array}\Bigr\}
\]
\note{La valeur de $\tau(\rho)$ est surchargée ($(\rho^{-1})$ biffé, puis un signe illisible) ; après $\sigma^{-1}$, un signe biffé.}
\emph{NB} $\tau =$ \struck{\ill{}} $\sigma\tau' = \tau'\sigma^{-1}$

\subsection{(I) Position relative de $L_\rho$ et $L_\sigma$}

\drawing{un diagramme en perspective d'inclusions de sous-groupes, chaque flèche à crochet marquée de son indice : $\{1\}\overset{2}{\hookrightarrow}\mu$, $\{1\}$ et $\mu_{\tau'}$ reliés par un trait d'indice 2, $\mu_{\tau'}\overset{2}{\hookrightarrow}D_{\tau'}$, $\mu$ et $D_{\tau'}$ reliés par un trait d'indice 2 ; $D_{\tau'}\overset{2}{\hookrightarrow}D_\sigma$, $D_{\tau'}\hookrightarrow\Pi_{\tau'}$, $D_{\tau'}\overset{3}{\hookrightarrow}D_\rho$ ; $\mu\overset{2}{\hookrightarrow}L_\sigma$, $\mu\hookrightarrow\underline{\Pi}$, $\mu\overset{3}{\hookrightarrow}L_\rho$ ; $L_\sigma$ et $\Pi_{\tau'}$ reliés par un trait d'indice 2, $L_\sigma$ et $\underline\Pi$ de même, $L_\rho$ et $D_\rho$ (indice 2) ; $\Pi_{\tau'}\overset{6}{\hookrightarrow}\mathfrak{S}$, $D_\sigma\hookrightarrow\mathfrak{S}$, $D_\rho\hookrightarrow\mathfrak{S}$, $D_\rho\overset{6}{\hookrightarrow}\underline{\mathfrak{S}}^+$, $L_\rho\hookrightarrow\underline{\mathfrak{S}}^+$, $\mathfrak{S}$ et $\mathfrak{S}^+$ reliés par un trait d'indice 2. Des étiquettes au crayon, $SD_\sigma$ et $SZ_\sigma$, sont esquissées.}

\drawing{à droite, le même diagramme passé au quotient par $\Pi_0$ : $\{1\}$ et $\mu$ (indice 2), $\mu$ et $D_2$ (indice 2), $D_2\overset{2}{\hookrightarrow}D_\sigma$, $L_\sigma$ et $D_\sigma$ (indice 2), $D_\sigma\overset{3}{\hookrightarrow}\mathfrak{S}/\Pi_0$, $D_2\overset{3}{\hookrightarrow}D_\rho$, $\mu\overset{2}{\hookrightarrow}L_\sigma$, $\mu\overset{3}{\hookrightarrow}L_\rho$, $L_\rho$ et $D_\rho$ (indice 2), $D_\rho\overset{2}{\hookrightarrow}\mathfrak{S}/\Pi_0$, $D_\rho\overset{3}{\hookrightarrow}\mathfrak{S}^+/\Pi_0$, $L_\rho\overset{2}{\hookrightarrow}\mathfrak{S}^+/\Pi_0$, $\mathfrak{S}/\Pi_0$ et $\mathfrak{S}^+/\Pi_0$ (indice 2).}

\[
  \text{ordre } 24\quad \mathfrak{S}/\Pi_0 = \Bigl\{\rho,\sigma,\tau \ \Bigm|\
  \begin{array}{l}
    \rho^6 = \sigma^4 = \tau^2 = 1,\ \rho^3 = \sigma^2\\
    \tau(\sigma) = \sigma^{-1},\ \tau(\rho) = \rho\\
    \sigma(\rho) = \rho^{-1}
  \end{array}\Bigr\}
\]
\[
  \text{ordre } 12\quad \mathfrak{S}^+/\Pi_0 = \Bigl\{\rho,\sigma \ \Bigm|\
  \begin{array}{l}
    \rho^6 = \sigma^4 = 1,\ \rho^3 = \sigma^2\\
    \sigma(\rho) = \rho^{-1}
  \end{array}\Bigr\}
\]

\subsection{(II)}

\drawing{un cube d'inclusions : face avant $\mu\times\mu_\tau = D_2 \overset{2}{\hookrightarrow} \mu\times\Pi'_0 = \mathfrak{S}^! = \Pi^!_0$ (lecture incertaine) $\overset{6}{\hookrightarrow}\mathfrak{S}$ ; $\mu_\tau\overset{2}{\hookrightarrow}\mathfrak{S}^!_0 = \Pi'$ ; face arrière $\underline\mu \overset{2}{\hookrightarrow} \mu\times\Pi_0 = \mathfrak{S}^{!+} = \Pi \overset{6}{\hookrightarrow} \underline{\mathfrak{S}}^+$, $\{1\}\overset{2}{\hookrightarrow}\mathfrak{S}^{!+}_0 = \Pi_0$ ; toutes les arêtes verticales et diagonales d'indice 2 ; $\mathfrak{S}$ et $\mathfrak{S}^+$ reliés par un trait d'indice 2. Une ligne courbe sépare ce cube de la partie (I).}

\[
  \underline\mu \simeq \mu_\tau \simeq \{\pm1\}
\]
\[
  \mu_\tau = \{1,\tau\},\quad \mu_{\tau'} = \{1,\tau'\}
\]
\[
  \mu = \{\pm1\}
\]
\[
  \mathfrak{S} =
\]
\[
  \mathfrak{S}^+ \qquad \rho\longmapsto\rho,\quad \sigma\longmapsto -\sigma
\]
\drawing{au crayon, très pâle, une esquisse d'un cube plus grand : $\mathcal{M}^!(0)\to\mathcal{M}^!_\Pi[0]$, $\mathcal{M}^!$, $\mu_\tau$, $D^\Pi_0$, $\mathcal{M}^!_{\mathfrak{S}}[0]$, $\mathcal{M}$, $1$, $L^\Pi_0$, $D^{\mathfrak{S}}_0$, $\Pi$, $\mathfrak{S}$, $L^{\mathfrak{S}}_0$, $\mathfrak{S}^+$, et plus bas $\Pi' = \mathfrak{S}^!$, $\Pi = \mathfrak{S}^{!+}$, $\mathfrak{S}^+$, $\mathcal{M}^{!+}$, $\mathcal{M}^+$ — premier état du cube de la p.~135.}

\page{135}
\struck{Position de} Les relations
\[
  \Pi \simeq \Pi_0\times\{\pm1\},\qquad
  \Pi' = \{1,\tau\}\cdot\Pi = D_\tau\cdot\Pi_0 = \Pi'_0\times\mu
\]
(sous $\Pi$ et $\Pi_0$ : $\wr\,\mathfrak{S}^{+!}$ et $\wr\,\mathfrak{S}^{+!}_0$)
\[
  \Pi'_0 = \{1,\tau\}\cdot\Pi_0
\]

(II)

\begin{tikzcd}[column sep=small, row sep=small, nodes={font=\scriptsize}]
\mu_\tau \arrow[rr, hook, "2"] \arrow[dd, no head, "2"'] \arrow[dr, hook, "2"] & & D_\tau = \mu_\tau\times\mu \arrow[dr, hook, "2"] \arrow[dd, no head, "2"'] & \\
 & \mathfrak{S}^!_0 = \Pi'_0 \arrow[rr, hook, "2"] \arrow[dd, no head, "2"'] & & \Pi' = \Pi'_0\times\mu = \mathfrak{S}^! \arrow[dd, no head, "2"] \\
\{1\} \arrow[rr, hook, "2"] \arrow[dr, hook, "2"] & & \underline\mu \arrow[dr, hook, "2"] & \\
 & \mathfrak{S}^{+!}_0 = \Pi_0 \arrow[rr, hook, "2"] & & \Pi = \Pi_0\times\mu = \mathfrak{S}^{+!}
\end{tikzcd}

\note{Les arêtes verticales sont tracées sans tête, avec un crochet à une extrémité : on les rend ici sans tête, marquées de leur indice 2.}

\begin{tikzcd}[column sep=small, row sep=small, nodes={font=\scriptsize}]
\mathfrak{S}^! \arrow[r, hook, "6"] \arrow[d, hook] & \mathfrak{S} \arrow[d, no head, "2"] \arrow[dr, hook] & \\
\mathcal{M}^! \arrow[rr, hook, "6"] \arrow[d, no head, "\hat{\mathbb{Z}}^*"'] & \mathfrak{S}^+ \arrow[dr, hook] & \mathcal{M} \arrow[d, no head, "\hat{\mathbb{Z}}^*"] \\
\mathcal{M}^{+!} \arrow[rr, hook, "6"] & & \mathcal{M}^+
\end{tikzcd}

\note{Dans l'original, $\mathfrak{S}^{+!} \overset{6}{\hookrightarrow} \mathfrak{S}^+$ (prolongeant la ligne $\Pi = \mathfrak{S}^{+!}$ du cube) et $\mathfrak{S}^{+!}\hookrightarrow\mathcal{M}^{+!}$ complètent ce cube ; les deux arêtes $\hat{\mathbb{Z}}^*$ portent chacune un point noir.}

\[
  \Pi_0 \simeq \{\lambda_0,\lambda_1,\lambda_\infty \mid \lambda_\infty\lambda_1\lambda_0 = 1\}
  \hookrightarrow \mathfrak{S},\qquad \lambda_i\longmapsto -l_i
\]
\drawing{une variante du cube précédent ($\Pi' = \mathfrak{S}^!$, $\mathcal{M}^!$, $\mathcal{M}$, $\mathcal{M}^{+!}$, $\mathcal{M}^+$, $\mathfrak{S}^{+!}\Rightarrow\mathfrak{S}^+$), couverte de boucles qui la cancellent.}

$\mathfrak{S}/\Pi_0$ ou $\mathfrak{S}^+/\Pi_0$ \uncertain{déduits} de $\mathfrak{S}$ \ill{} $\mathfrak{S}^+$ en divisant
par la relation $\tau'(\rho) =$ \struck{\ill{}} $\sigma(\rho)$ ou encore $\tau(\rho) = \rho$
(\emph{NB} $\tau(\rho) =$ \uncertain{$1-(l_1^{-1})\rho$} !)

\begin{tikzcd}[column sep=small, row sep=small, nodes={font=\scriptsize}]
\mathcal{M}(0) \arrow[r, leftrightarrow] \arrow[d, no head] & \mathcal{M}^{\Pi}[0] \arrow[rr, hook] \arrow[d, no head] \arrow[dr, hook, "6"] & & \mathcal{M}^! \arrow[dr, hook] \arrow[d, no head] & \\
\mu_\tau \arrow[r, hook] \arrow[d, no head] & D^{\Pi}_0 \arrow[d, no head] \arrow[dr, hook, "2"] & \mathcal{M}^{\mathfrak{S}}[0] \arrow[d, no head] & \Pi' = \mathfrak{S}^! \arrow[dr, hook] \arrow[d, no head] & \mathcal{M} \arrow[d, no head] \\
\{1\} \arrow[r, hook, "\hat{\mathbb{Z}}"] & L^{\Pi}_0 \arrow[dr, hook, "2"] & D^{\mathfrak{S}}_0 \arrow[d, no head] & \underline\Pi = \mathfrak{S}^{+!} \arrow[dr, hook, "6"] & \mathfrak{S} \arrow[d, no head] \\
 & & L^{\mathfrak{S}}_0 & & \mathfrak{S}^+
\end{tikzcd}

\note{Cube redessiné de façon simplifiée : sur la page, $D^{\Pi}_0\hookrightarrow\Pi'$, $L^{\Pi}_0\hookrightarrow\underline\Pi$, $\mathcal{M}^{\mathfrak{S}}[0]\hookrightarrow\mathcal{M}$, $D^{\mathfrak{S}}_0\hookrightarrow\mathfrak{S}$ et $L^{\mathfrak{S}}_0\hookrightarrow\mathfrak{S}^+$ sont des flèches horizontales qui traversent les arêtes verticales ; l'exposant de $\mathcal{M}^{\Pi}[0]$ est surchargé. Au crayon, sous le cube, $L^{\Pi_0}_0$.}

\subsection{III}

\struck{position de $\mathcal{M}(0)$}

$\mathcal{M}(0)$ et groupes associés

\page{137}
III \quad $\mathcal{M}(0)$

\drawing{un grand diagramme en perspective (trois couches de cubes accolés) d'inclusions marquées de leurs indices. Couche supérieure : $\mathcal{M}(0)\overset{2}{\hookrightarrow}\mathcal{M}(0)\times\mu$, $\mathcal{M}^!_0\overset{2}{\hookrightarrow}\underline{\mathcal{M}}^! \simeq \mathcal{M}^!_0\times\mu \overset{6}{\hookrightarrow}\mathcal{M}$, $\mathcal{M}^{\Pi_0}_{[0]}$, $\mathcal{M}^{\Pi}[0]$ (ou $\mathcal{M}^![0]$) $\overset{2}{\hookrightarrow}\mathcal{M}^{\mathfrak{S}}[0]$ ou $\mathcal{M}[0]$. Couche médiane : $\mu_\tau$, $D_\tau$, $\{1\}\overset{2}{\hookrightarrow}\Pi'_0 = \mathfrak{S}^!_0\overset{2}{\hookrightarrow}\Pi' = \mathfrak{S}^!\overset{6}{\hookrightarrow}\mathfrak{S}$, $\underline\mu$, $D^{\Pi_0}_0\overset{2}{\hookrightarrow}D^{\Pi}_0\overset{2}{\hookrightarrow}D^{\mathfrak{S}}_0$. Couche inférieure : $\underline\Pi_0\overset{2}{\hookrightarrow}\underline\Pi = \mathfrak{S}^{+!}\overset{6}{\hookrightarrow}\underline{\mathfrak{S}}^+$, $L^{\Pi_0}_0\overset{2}{\hookrightarrow}L^{\Pi}_0\overset{2}{\hookrightarrow}L^{\mathfrak{S}}_0$. Arêtes verticales et diagonales d'indice 2 ; à gauche, des arcs pointillés relient $\mathcal{M}(0)$, $\mu_\tau$, $\{1\}$ aux sommets $\mathcal{M}^{\Pi_0}_{[0]}$, $D^{\Pi_0}_0$, $L^{\Pi_0}_0$.}

IV \quad $Z_\rho,\ N_\rho$

\drawing{un second diagramme en perspective, même disposition. On y lit : $\mu_{\tau'}\cdot Z^!_{0\rho} = N^{?}_{0\rho}$ (biffé) $\hookrightarrow N^{?}_\rho$, avec au-dessus « $\mu_{\tau'}\cdot Z^!_0$ \ill{} » en partie biffé ; $Z^!_\rho\overset{6}{\hookrightarrow}N_\rho\hookrightarrow\mathcal{M}$ ; $Z^!_{0\rho}\overset{2}{\hookrightarrow}Z^!_\rho\overset{3}{\hookrightarrow}Z_\rho$ ; $\mathcal{M}^!$, $N_\rho$ (indice 6), $\mathcal{M}'$ ; $\mu_{\tau'}\hookrightarrow D_{\tau'}\hookrightarrow D_\rho$ ; $\mathcal{M}'^!_0\overset{2}{\hookrightarrow}\mathcal{M}'^!\overset{6}{\hookrightarrow}\mathcal{M}'$ ; $1\overset{2}{\hookrightarrow}\mu\overset{3}{\hookrightarrow}L_\rho$ ; $\underline\Pi_0\overset{2}{\hookrightarrow}\underline\Pi = \mathfrak{S}^{+!}\overset{6}{\hookrightarrow}\underline{\mathfrak{S}}^+$ ; arêtes verticales et diagonales d'indices 2 et 3.}

\page{139}
IV

\drawing{reprise au propre du diagramme IV de la p.~137, en trois blocs de cubes. Ligne supérieure : $N^?_{0\rho}$ (biffé) $\overset{2}{\hookrightarrow} N^?_\rho$ ; $N^?_\rho$ (biffé) ; $N^?_\rho \hookrightarrow \mathcal{M}^? \overset{6}{\hookrightarrow} \mathcal{M}$ ; $\mathcal{M}^?_0\overset{2}{\hookrightarrow}\mathcal{M}^!$. Ligne médiane : $Z^!_{0\rho} = \mathcal{M}(\bar{\jmath})\overset{2}{\hookrightarrow}Z^!_\rho\hookrightarrow\mathcal{M}'^!\overset{6}{\hookrightarrow}\mathcal{M}'$ ; $\mu_{\tau'}\overset{2}{\hookrightarrow}D_{\tau'}\hookrightarrow\Pi_{\tau'}\overset{6}{\hookrightarrow}\mathfrak{S}$ ; $\mathcal{M}'^!_0\hookrightarrow\mathcal{M}'^!$ ; $1\overset{2}{\hookrightarrow}\underline\mu\hookrightarrow\underline\Pi\overset{6}{\hookrightarrow}\underline{\mathfrak{S}}^+$. Ligne inférieure : $\underline\Pi_0\overset{2}{\hookrightarrow}\underline\Pi$, $\mathcal{M}'$ ; puis $\underline{\mathfrak{S}}^+$ au bas d'une arête d'indice 6. Arêtes verticales et diagonales d'indice 2 ; deux arcs pointillés relient $\underline\Pi$ et $\underline{\mathfrak{S}}^+$ (bloc de droite) à $\underline\Pi$ et $\underline{\mathfrak{S}}^+$ (bloc de gauche). Au crayon, en bas de page, une esquisse antérieure du même diagramme : $N^?_{0\rho}$ (biffé), $N^?_\rho$, $N_\rho$, $\mathcal{M}$ ; $Z^!_{0\rho} = \mathcal{M}(\bar{\jmath})$, $Z^!_\rho$, $Z_\rho$, $\mathcal{M}'$ ; $\mu_{\tau'}$, $D_{\tau'}$, $D_\rho$, $\mathfrak{S}$ ; $\mathcal{M}'^!_0$, $\mathcal{M}'^!$ ; $\underline\mu$, $L_\rho$, $\underline{\mathfrak{S}}^+$ ; $\underline\Pi_0$, $\underline\Pi$.}
\note{Les étiquettes lues « 6 » sur les arêtes horizontales ont ici la forme d'un $\sigma$ (lecture incertaine) ; sur la p.~137 les arêtes correspondantes portent nettement « 6 ».}

\section{$S\mathcal{T}_{0,2}$, $S\mathcal{T}_{1,1}$, $S\mathcal{T}_{1,2}$ …}
\note{Titre inscrit en haut d'une feuille (p.~140) restée blanche par ailleurs, une couverture : ce sont ses mots. Ce qu'elle annonce se trouve vraisemblablement dans le lot suivant. Dans $S\mathcal{T}_{1,1}$ le second indice est surchargé. Le dernier mot, rendu « … », se lit peut-être \uncertain{etc}.}

\end{document}
